% answers_short.tex —— 2009 年·普通高等学校招生考试·浙江卷·理科 · 简短答案速查
% 由 tools/make_short_answers.py 自 answers.tex 的 \gksolhead/\gkanswer 提取，
% 并与 2009/2009_数学真题_浙江理_答案要点.md 逐题比对后生成，勿手改（改源 MD 后重跑）。

\begin{gknotice}
\textbf{使用说明}：本节只给最终答案便于快速核对，完整推导见解析版。答案为编写者独立推导，\textbf{非官方标准答案}。
\end{gknotice}

\gkpart{选择题}
\begin{center}\begin{tabular}{ccccccccccc}
\toprule
题号 & 1 & 2 & 3 & 4 & 5 & 6 & 7 & 8 & 9 & 10 \\
\midrule
答案 & B & C & D & B & C & A & B & D & C & C \\
\bottomrule
\end{tabular}\end{center}

\gkpart{填空题}
\begin{itemize}[leftmargin=2.6em]
  \item[\textbf{11.}] $15$。
  \item[\textbf{12.}] $18$。
  \item[\textbf{13.}] $4$。
  \item[\textbf{14.}] $148.4$。
  \item[\textbf{15.}] $2^{4n-1}+(-1)^n2^{2n-1}$。
  \item[\textbf{16.}] $336$。
  \item[\textbf{17.}] $\left(\dfrac12,1\right)$。
\end{itemize}

\gkpart{解答题}
\begin{itemize}[leftmargin=2.6em]
  \item[\textbf{18.}] （1）$2$；（2）$a=2\sqrt5$。
  \item[\textbf{19.}] （1）$\dfrac{10}{21}$；（2）$P(\xi=0)=\dfrac{5}{12}$，$P(\xi=1)=\dfrac12$，$P(\xi=2)=\dfrac1{12}$，$E\xi=\dfrac23$。
  \item[\textbf{20.}] （1）证明见解析；（2）$M$ 存在，$M$ 到 $OA$ 的距离为 $4$，到 $OB$ 的距离为 $\dfrac94$。
  \item[\textbf{21.}] （1）$\dfrac{y^2}{4}+x^2=1$；（2）$h$ 的最小值为 $1$。
  \item[\textbf{22.}] （1）$-5<k<-2$；（2）存在，$k=5$。
\end{itemize}
